\spacing{1.3}


\section{Material Systems}
We can classify material systems into (see diagram):
\vspace{12pt}

\esquema{Material Systems}
{\esquema{Pure Substances}
{\esquema{Elements}
{\( \bullet \) Hydrogen
\vspace{6pt}\\
    \( \bullet \) Lithium}
\vspace{20pt}\\
\esquema{Compounds}
{\( \bullet \)\ch{NaCl}
\vspace{6pt}\\
    \( \bullet \)\ch{H2O}
\vspace{6pt}\\
    \( \bullet \)\ch{H2SO4}}
}
\vspace{20pt} \\
\esquema{Mixtures}
{\( \bullet \) Homogeneous
\vspace{6pt}\\
    \( \bullet \) Heterogeneous
}
}

\begin{enumerate}[label=\numberballmDarkTeal]
\item \negsubro{Pure substances}: a type of matter that has a constant composition and defined properties, and that cannot be separated into other substances by physical methods.
Pure substances can be divided into:
\begin{itemize}
    \item \textbf{Elements}: made of only one type of atom.
    \item \textbf{Compounds}: made by the union of two or more chemical elements.
\end{itemize}

\item \negsubro{Mixtures}: the combination of two or more pure substances that can be separated by physical methods.
Mixtures can be divided into:
\begin{itemize}
    \item \textbf{Homogeneous mixture}: a mixture where the components are evenly distributed and cannot be seen with the naked eye.
    \item \textbf{Heterogeneous mixture}: a mixture where the components can be distinguished.
\end{itemize}

\end{enumerate}


\section{Solutions}

\begin{definition}[Solution]
    Is a homogeneous system made of two or more pure substances in different proportions.
\end{definition}

The components of a solution are:
\begin{itemize}
    \item \textbf{Solvent}. The substance in greater amount.
    \item \textbf{Solute}. The substance in smaller amount. There can be more than one solute.
    \item \textbf{Solution}. The mixture of solute(s) and solvent.
\end{itemize}

For example, in a solution of water and sugar, the solvent is water and the solute is sugar.
In a bottle of mineral water, the label says it contains sodium, calcium, magnesium, etc.
In this case, the solvent is water, and the solutes are sodium, calcium, and magnesium.

\subsection{Concentration of a solution}

When we have a solution, it is crucial to know how much solute there is in the solvent.

\begin{definition}[Concentration]
    Is the amount of solute dissolved in a certain amount of solvent.
\end{definition}

There are different ways to express concentration. The main ones are:

\begin{enumerate}[label=\numberballmDarkTeal]
\item \negsubro{Mass percentage}.
It tells us the mass of solute in every \qty{100}{\g} of solution.

\begin{equation}
    \text{Mass \% of solute} = \dfrac{\text{mass of solute}}{\text{mass of solution}}\cdot 100
    \label{eq:porcentaje_masa}
\end{equation}

Important:
\begin{itemize}
    \item If there are several solutes, we calculate the percentage of each one.
    \item We can also calculate the percentage of solvent.
    \item All percentages together must add up to \num{100}.
\end{itemize}
\end{enumerate}

\begin{example}
    We have a water solution with sodium and calcium.
    There are \qty{18}{\g} of sodium, \qty{6}{\g} of calcium, and \qty{100}{\g} of water.
    Calculate the mass percentage of sodium and calcium.

    For both cases we must apply equation \eqref{eq:porcentaje_masa}.
    In both cases we have to:

    \begin{equation*}
        \text{\% mass of sodium} = \dfrac{\qty{18}{\g}\text{ of sodium}}{\qty{18}{\g}\text{ of sodium} + \qty{6}{\g}\text{ of calcium} + \qty{100}{\g}\text{ of water}}\cdot 100 = 14.52 \%
    \end{equation*}

    \begin{equation*}
        \text{\% mass of calcium} = \dfrac{\qty{6}{\g}\text{ of calcium}}{\qty{18}{\g}\text{ of sodium} + \qty{6}{\g}\text{ of calcium} + \qty{100}{\g}\text{ of water}}\cdot 100 = 4.84 \%
    \end{equation*}

    If we also want to calculate the percentage of solvent, it would be:

    \begin{equation*}
        \text{\% mass of water} = \dfrac{\qty{100}{\g}\text{ of water}}{\qty{18}{\g}\text{ of sodium} + \qty{6}{\g}\text{ of calcium} + \qty{100}{\g}\text{ of water}}\cdot 100 = 80.65 \%
    \end{equation*}
\end{example}

\begin{exercise}
    Steel is an alloy composed of iron, carbon and other elements. A \qty{120}{\g} mass of steel contains 1.8\% carbon and 1.2\% manganese. Find the mass of each element.
    \begin{solu}
        \textbf{Formula:} $m_{\text{component}} = \frac{\% \text{ by mass}}{100} \cdot m_{\text{total}}$ \\
        \textbf{Data:} $m_{\text{total}} = 120 \text{ g}$; $\%C = 1.8\%$; $\%Mn = 1.2\%$ \\
        \textbf{Solution:}
        \begin{enumerate}
            \item \textbf{Mass of carbon:} $m_{C} = \frac{1.8}{100} \cdot 120 \text{ g} = 2.16 \text{ g}$
            \item \textbf{Mass of manganese:} $m_{Mn} = \frac{1.2}{100} \cdot 120 \text{ g} = 1.44 \text{ g}$
            \item \textbf{Mass of iron:} It is calculated by subtracting the other components from the total: \\
            $m_{Fe} = 120 \text{ g} - (2.16 \text{ g} + 1.44 \text{ g}) = 116.4 \text{ g}$
        \end{enumerate}
    \end{solu}
\end{exercise}

\begin{exercise}
    We dissolve 12 g of sodium chloride and 13 g of potassium chloride in 250 g of distilled water. Find the percentage by mass of each solute in the final solution.
    \begin{solu}
        \textbf{Formulas:} $m_{\text{total}} = m_{s1} + m_{s2} + m_{\text{solvent}}$ and $\%m = \frac{m_{\text{solute}}}{m_{\text{total}}} \cdot 100$ \\
        \textbf{Solution:}
        \begin{enumerate}
            \item \textbf{Total mass:} $12 \text{ g} + 13 \text{ g} + 250 \text{ g} = 275 \text{ g}$.
            \item \textbf{\% NaCl:} $\frac{12 \text{ g}}{275 \text{ g}} \cdot 100 \approx 4.36\%$
            \item \textbf{\% KCl:} $\frac{13 \text{ g}}{275 \text{ g}} \cdot 100 \approx 4.73\%$
        \end{enumerate}
    \end{solu}
\end{exercise}

\begin{exercise}
    We want to prepare 500 g of an aqueous solution containing 0.5\% sodium chloride, 0.3\% potassium chloride and 0.2\% potassium nitrate by mass. How much of each solute is needed? How much water do we need?
    \begin{solu}
        \textbf{Formula:} $m_s = \frac{\%m}{100} \cdot m_{\text{total}}$ \\

        \textbf{Data:} $m_{\text{total}} = 500 \text{ g}$. \\

        \textbf{Solution:}
        \begin{itemize}
            \item $m_{NaCl} = 500 \cdot 0.005 = 2.5 \text{ g}$
            \item $m_{KCl} = 500 \cdot 0.003 = 1.5 \text{ g}$
            \item $m_{KNO_3} = 500 \cdot 0.002 = 1.0 \text{ g}$
            \item \textbf{Mass of water:} $500 \text{ g} - (2.5 + 1.5 + 1.0) \text{ g} = 495 \text{ g}$
        \end{itemize}
    \end{solu}
\end{exercise}

\begin{exercise}
    What does it mean that a solution of iodine in alcohol is 1\% by mass?
    \begin{solu}
        \textbf{Definition of mass percentage:} It represents the parts of solute in 100 parts of solution. \\
        It means that in every \textbf{100 g of solution}, there is exactly \textbf{1 g of iodine} (solute) and the remaining \textbf{99 g are alcohol} (solvent).
    \end{solu}
\end{exercise}

\begin{exercise}
    We use a mixture of lead and tin in soldering. If the percentage by mass of lead is 66\% and that of tin is 34\%, what masses of each metal would you melt to prepare 250 g of this alloy?
    \begin{solu}
        \textbf{Formula:} $m_{\text{component}} = \frac{\%m}{100} \cdot m_{\text{total}}$ \\

        \textbf{Data:} $m_{\text{total}} = 250 \text{ g}$. \\

        \textbf{Solution:}
        \begin{itemize}
            \item Lead mass (Pb): $250 \cdot 0.66 = 165 \text{ g}$
            \item Tin mass (Sn): $250 \cdot 0.34 = 85 \text{ g}$
        \end{itemize}
        (Check: $165 + 85 = 250 \text{ g}$).
    \end{solu}
\end{exercise}


\begin{exercise}
    The percentage by mass of potassium chloride in a certain solution is 10\%. How much of this solute is dissolved in 250 g of solution?
    \begin{solu}
        \textbf{Formula:} $m_s = \frac{\%m}{100} \cdot m_{\text{solution}}$ \\
        \textbf{Solution:} We apply the definition directly to the total mass: \\
        $m_{KCl} = 250 \text{ g} \cdot \frac{10}{100} = 25 \text{ g}$ of KCl.
    \end{solu}
\end{exercise}

\begin{exercise}
    A bottled water contains 99.8\% by mass of water and 0.2\% carbon dioxide. Identify the solute and solvent in this solution. What mass of carbon dioxide would 1 kg of water contain?
    \begin{solu}
        \textbf{Concept:} Identify solute (smaller amount) and solvent (greater amount). \\
        \textbf{Proportion formula:} $\frac{m_{\text{solute}}}{m_{\text{solvent}}} = \frac{\% \text{ solute}}{\% \text{ solvent}}$ \\
        \textbf{Solution:}
        \begin{itemize}
            \item \textbf{Identification:} Solute = carbon dioxide ($CO_2$); Solvent = water ($H_2O$).
            \item \textbf{Mass calculation:} We have 1 kg (1000 g) of water. \\
            $m_{CO_2} = \frac{0.2 \text{ g } CO_2}{99.8 \text{ g } H_2O} \cdot 1000 \text{ g } H_2O \approx 2.004 \text{ g}$ of $CO_2$.
        \end{itemize}
    \end{solu}
\end{exercise}

\begin{enumerate}[label=\numberballmDarkTeal,resume]
\item \negsubro{Volume percentage}.
It tells us the volume of solute in every \qty{100}{\litre} of solution.

\begin{equation}
    \text{Volume \% of solute} = \dfrac{\text{volume of solute}}{\text{volume of solution}}\cdot 100
\end{equation}
\end{enumerate}

\begin{example}
    We prepare a solution with \qty{100}{\ml} of alcohol, \qty{150}{\ml} of acetone, and \qty{250}{\ml} of water.
    Find the volume percentage of each substance.

    \begin{equation*}
        \text{\% vol alcohol} = \dfrac{\qty{100}{\ml}\text{ of alcohol}}{\qty{100}{\ml}\text{ of alcohol} + \qty{150}{\ml}\text{ of acetone} + \qty{250}{\ml}\text{ of water}}\cdot 100 = 20 \%
    \end{equation*}

    \begin{equation*}
        \text{\% vol acetone} = \dfrac{\qty{150}{\ml}\text{ of acetone}}{\qty{100}{\ml}\text{ of alcohol} + \qty{150}{\ml}\text{ of acetone} + \qty{250}{\ml}\text{ of water}}\cdot 100 = 30 \%
    \end{equation*}

    \begin{equation*}
        \text{\% vol water} = \dfrac{\qty{250}{\ml}\text{ of water}}{\qty{100}{\ml}\text{ of alcohol} + \qty{150}{\ml}\text{ of acetone} + \qty{250}{\ml}\text{ of water}}\cdot 100 = 50 \%
    \end{equation*}
\end{example}

\begin{exercise}
    A solution of ether and chloroform is prepared by adding \qty{10}{\mL} of ether to 90 mL of chloroform. What is the percentage by volume of each component?
    \begin{solu}
        \textbf{Formulas:} \\
        1. $V_{\text{solution}} = V_{\text{solute}} + V_{\text{solvent}}$ \\
        2. $\%V = \frac{V_{\text{component}}}{V_{\text{solution}}} \cdot 100$ \\

        \textbf{Data:} $V_{\text{ether}} = 10 \text{ mL}$; $V_{\text{chloroform}} = 90 \text{ mL}$ \\

        \textbf{Solution:}
        \begin{enumerate}
            \item \textbf{Volume of the solution:} $V_{\text{total}} = 10 \text{ mL} + 90 \text{ mL} = 100 \text{ mL}$.
            \item \textbf{Ether percentage:} $\%V_{\text{ether}} = \frac{10 \text{ mL}}{100 \text{ mL}} \cdot 100 = 10\%$
            \item \textbf{Chloroform percentage:} $\%V_{\text{chloro}} = \frac{90 \text{ mL}}{100 \text{ mL}} \cdot 100 = 90\%$
        \end{enumerate}
    \end{solu}
\end{exercise}

\begin{exercise}
    The composition by volume of the active ingredients in a 200 mL bottle of body lotion is: 2.14\% sweet almond oil; 2.14\% corn oil; 4.29\% peanut oil and 2.88\% hydrogenated copra oil. Calculate the volume of each component.
    \begin{solu}
        \textbf{Formula:} $V_{\text{comp}} = \frac{\%V}{100} \cdot V_{\text{total}}$ \\
        \textbf{Data:} $V_{\text{total}} = 200 \text{ mL}$. \\
        \textbf{Solution:} We multiply the total volume by each decimal percentage:
        \begin{itemize}
            \item Almond oil: $200 \cdot 0.0214 = 4.28 \text{ mL}$
            \item Corn oil: $200 \cdot 0.0214 = 4.28 \text{ mL}$
            \item Peanut oil: $200 \cdot 0.0429 = 8.58 \text{ mL}$
            \item Copra oil: $200 \cdot 0.0288 = 5.76 \text{ mL}$
        \end{itemize}
    \end{solu}
\end{exercise}


\begin{exercise}
    We have dissolved \qty{25}{\cmcubicos} of alcohol in distilled water to obtain a total solution volume of \qty{250}{\cmcubicos}.
    \begin{enumerate}[label=\letraminuscula]
        \item What is the volume of the solute and the solvent?
        \item Calculate the percentage by volume of each component of this solution.
    \end{enumerate}
    \begin{solu}
        \textbf{Formulas:} $V_{\text{solvent}} = V_{\text{total}} - V_{\text{solute}}$ and $\%V = \frac{V_{\text{comp}}}{V_{\text{total}}} \cdot 100$ \\
        \textbf{Solution:}
        \begin{enumerate}[label=\alph*)]
            \item \textbf{Volumes:} Solute = $25 \text{ cm}^3$ (alcohol). \\
            Solvent = $250 \text{ cm}^3 - 25 \text{ cm}^3 = 225 \text{ cm}^3$ (water).
            \item \textbf{Percentages:} \\
            $\%V_{\text{alcohol}} = \frac{25}{250} \cdot 100 = 10\%$ \\
            $\%V_{\text{water}} = \frac{225}{250} \cdot 100 = 90\%$
        \end{enumerate}
    \end{solu}
\end{exercise}


\begin{enumerate}[label=\numberballmDarkTeal,resume]
\item \negsubro{Mass concentration}.
It tells us the mass of solute per unit of volume of solution.
It is usually expressed in \unit{\g\per\litre}.

\begin{equation*}
    \text{Mass concentration} = \dfrac{m_{\text{solute (g)}}}{V_{\text{solution (L)}}}
\end{equation*}

Do not confuse it with density:
\begin{itemize}
    \item Density is the mass of a substance per unit of volume.
    \item Mass concentration is the mass of solute per unit of volume of the \textbf{solution}.
\end{itemize}
\end{enumerate}

\begin{example}
    To prepare a potassium nitrate solution with \qty{20}{\g} of nitrate in water to obtain a final solution volume of \qty{200}{\ml}. Calculate:
    \begin{enumerate}[label=\letraminuscula]
\item The mass concentration of the solution expressed in \unit{\g\per\litre}.

Since the volume of the solution is given in \unit{\ml}, we convert it to \unit{\litre}. The volume is then \qty{0.2}{\litre}. If we apply:
\begin{equation*}
    \text{Mass concentration} = \dfrac{m_{\text{solute (g)}}}{V_{\text{solution (L)}}} = \dfrac{20}{0.2} = \qty{100}{\g\per\litre}
\end{equation*}

\item Calculate the mass of the solution if the density of this solution is \( \rho = \qty{1.09}{\g\per\cubic\centi\metre} \).

To calculate its mass, we solve the density equation for mass:
\begin{equation*}
    \rho = \dfrac{m}{V} \rightarrow m = \rho \cdot V
\end{equation*}
Before substituting into the equation, we must check that both the density and the volume are in the same units, but we know that \( \qty{1}{\ml} = \qty{1}{\cubic\centi\metre} \).
\begin{equation*}
    m = \rho \cdot V = 1.09 \cdot 200 = \qty{218}{\g}
\end{equation*}
    \end{enumerate}
\end{example}

\begin{exercise}
    200 ml of an aqueous solution contains 12 g of sugar. The density of this solution is \qty{1.022}{\g\per\cmcubicos}.
    \begin{enumerate}[label=\letraminuscula]
        \item Calculate the percentage by mass of the sugar.
        \item Find the mass concentration in g/L.
    \end{enumerate}
    \begin{solu}
        \textbf{Formulas:} $m_{\text{solution}} = d \cdot V$; $\%m = \frac{m_s}{m_{\text{solution}}} \cdot 100$; $C = \frac{m_s}{V(L)}$ \\
        \textbf{Solution:}
        \begin{enumerate}[label=\alph*)]
        \item \textbf{Mass percentage:}
        First we find the total mass: $m_{\text{total}} = 1.022 \text{ g/cm}^3 \cdot 200 \text{ cm}^3 = 204.4 \text{ g}$. \\
        $\%m = \frac{12 \text{ g}}{204.4 \text{ g}} \cdot 100 \approx 5.87\%$
        \item \textbf{Concentration in g/L:}
        We divide the solute mass by the volume in litres ($200 \text{ mL} = 0.2 \text{ L}$): \\
        $C = \frac{12 \text{ g}}{0.2 \text{ L}} = 60 \text{ g/L}$
        \end{enumerate}
    \end{solu}
\end{exercise}

\begin{exercise}
    An isotonic saline infusion can be prepared by dissolving 0.9 g of sodium chloride in water to give 100 mL of solution.
    Calculate:
    \begin{enumerate}[label=\letraminuscula]
            \item the mass concentration in g/L of this solution.
            \item the mass of sodium chloride needed to prepare 250 mL of isotonic saline.
    \end{enumerate}
    \begin{solu}
        \textbf{Formulas:} $C \text{ (g/L)} = \frac{m_{\text{solute (g)}}}{V_{\text{solution (L)}}}$ and $m_s = C \cdot V$ \\
        \textbf{Solution:}
        \begin{enumerate}[label=\alph*)]
                \item \textbf{Concentration:} We change the volume to litres ($100 \text{ mL} = 0.1 \text{ L}$). \\
                $C = \frac{0.9 \text{ g}}{0.1 \text{ L}} = 9 \text{ g/L}$.
                \item \textbf{Mass for 250 mL:} We change the volume ($250 \text{ mL} = 0.25 \text{ L}$). \\
                $m = 9 \text{ g/L} \cdot 0.25 \text{ L} = 2.25 \text{ g}$ of NaCl.
        \end{enumerate}
    \end{solu}
\end{exercise}


\section{Solubility}

\begin{definition}[Solubility]
        The \textbf{solubility} of a solute at a certain temperature and in a certain solvent is the maximum amount of that solute that can dissolve in \qty{100}{\g} of solvent.
\end{definition}

From this definition we can see that solubility depends on temperature and on the solvent.
This means that solubility changes when temperature changes and when the solvent changes.

For example, if something does not dissolve at room temperature, we can heat it.
This shows that when the temperature increases, solubility also increases.

To express solubility correctly, we must give the value of the solubility, say that it is in \qty{100}{\g} of the solvent used, and also give the temperature.

\subsection{Solubility curves}\label{subsec:curvas-de-solubilidad}

A \textbf{solubility curve} is a graph that shows the solubility data of a pure substance in \qty{100}{\g} of a solvent at different temperatures.

The Figure \ref{fig:curva_solubilidad}, shows the solubility curves of different pure substances in water are shown.

\section*{Summary}
\begin{itemize}
    \item \textbf{Solute:} The substance that dissolves and is usually present in a smaller amount in the mixture.
    \item \textbf{Solvent:} The substance that dissolves the solute. It is usually the main component (for example, water in aqueous solutions).
    \item \textbf{Solution:} The homogeneous mixture formed by the solute and the solvent ($m_{\text{solution}} = m_{\text{solute}} + m_{\text{solvent}}$) or, for volume, \( V_{\text{solution}} = V_{\text{solute}} + V_{\text{solvent}} \).
    \item \textbf{Concentration:} It expresses the proportion of solute in a certain amount of solution or solvent.
    \item \textbf{Density ($d$):} Relationship between the total mass of the solution and the volume it occupies ($d = m/V$). It is used to change volume into mass and mass into volume. It is important not to confuse density with mass/volume concentration.
\end{itemize}

\begin{exercise}
    Use the solubility curves on the last page to answer these questions.
    \begin{enumerate}[label=\letraminuscula]
            \item What’s the solubility of KNO3 at 25 °C and 45 °C?
            \item What mass of \ch{KNO3} crystals will form if a saturated solution of 100 g of water is cooled from 45 °C to 25 °C? And from 50 °C to 20 °C?
            \item What’s the greatest mass of KNO3 that could dissolve in 1 kg of water at 50 °C?
            \item How much water is needed to dissolve 100 g of potassium nitrate at 45 °C?
            \item At what temperature does potassium nitrate have a solubility of 20 g per 100 g of water?
    \end{enumerate}
    \begin{solu}
        Aquí va la solución sería otra solución y mejor hacerlo todo.
    \end{solu}
\end{exercise}

\begin{exercise}
    Use the solubility curves of \ch{KCl} and \ch{NaOH} on the last page to answer the questions.
    \begin{enumerate}[label=\letraminuscula]
            \item The solubility of which of these two substances varies less with temperature?
            \item What are their solubilities at \qty{10}{\celsius}?
            \item What’s the solubility of \ch{KCl} at \qty{30}{\celsius}? And of \ch{NaOH} at 80 °C?
            \item What mass of \ch{KCl} will form if a solution saturated in 100 g of water is cooled from
            30 °C to 10 °C?
            \item What mass of \ch{NaOH} crystals will form if a solution saturated in 100 g of water is cooled
            from 80 °C to 10 °C?
            \item At what temperature does sodium chlorate have solubility of 90 g in 100 g of water?
    \end{enumerate}
    \begin{solu}
        Aquí va la solución
    \end{solu}
\end{exercise}

\section{Vocabulary}


\begin{vocabulary}[Matter]
Anything that has mass and occupies space. \newline
\textbf{Translation to Spanish}: Materia \newline
% Line 14
\end{vocabulary}



\begin{vocabulary}[Atom]
The smallest unit of an element that retains the chemical properties of that element. \newline
\textbf{Translation to Spanish}: Átomo \newline
% Line 41
\end{vocabulary}


\begin{vocabulary}[Molecule]
A group of atoms chemically bonded together. \newline
\textbf{Translation to Spanish}: Molécula \newline
% Line 57
\end{vocabulary}


\begin{vocabulary}[Element]
A pure substance consisting of only one type of atom. \newline
\textbf{Translation to Spanish}: Elemento \newline
% Line 76
\end{vocabulary}


\begin{vocabulary}[Compound]
A substance formed when two or more chemical elements are chemically combined. \newline
\textbf{Translation to Spanish}: Compuesto \newline
% Line 94
\end{vocabulary}


\begin{vocabulary}[Mixture]
A combination of substances that are not chemically bonded and can be separated physically. \newline
\textbf{Translation to Spanish}: Mezcla \newline
% Line 128
\end{vocabulary}


\begin{vocabulary}[Density]
The mass per unit volume of a substance. \newline
\textbf{Translation to Spanish}: Densidad \newline
% Line 161
\end{vocabulary}


\begin{vocabulary}[Solution]
A homogeneous mixture composed of a solute dissolved in a solvent. \newline
\textbf{Translation to Spanish}: Disolución \newline
% Line 219
\end{vocabulary}


\begin{vocabulary}[Solute]
The substance that is dissolved in a solution. \newline
\textbf{Translation to Spanish}: Soluto \newline
% Line 231
\end{vocabulary}


\begin{vocabulary}[Solvent]
The substance in which the solute dissolves to form a solution. \newline
\textbf{Translation to Spanish}: Disolvente \newline
% Line 235
\end{vocabulary}


\begin{vocabulary}[Chemical Change]
A process in which substances are transformed into new substances with different properties. \newline
\textbf{Translation to Spanish}: Cambio químico \newline
% Line 307
\end{vocabulary}


\begin{vocabulary}[Physical Change]
A change that affects the form or state of matter without altering its chemical composition. \newline
\textbf{Translation to Spanish}: Cambio físico \newline
% Line 329
\end{vocabulary}


\begin{vocabulary}[State of Matter]
One of the distinct forms that matter can exist in, such as solid, liquid, gas, or plasma. \newline
\textbf{Translation to Spanish}: Estado de la materia \newline
% Line 371
\end{vocabulary}



\newpage

\begin{figure}[H]
    \centering
    \begin{tikzpicture}
        \begin{axis}
            [
% --- Control del Grid ---
            grid=both,
            major grid style={line width=.2pt,draw=gray!70},
            minor grid style={line width=.1pt,draw=gray!50},
% ------------------------
            width=0.9\textwidth,
            height=0.7\textwidth,
            axis x line=bottom,
            axis y line=left,
            axis line style=thick,
            xmax=160, xmin=0,
            ymax=160, ymin=0,
            xtick={0,20,40,60,80,100, 120, 140, 160},
            ytick={0,20,40,60,80,100,120,140,160},
            minor tick num=3, % Esto crea las líneas pequeñas del grid
            xlabel={Temperatura ( \unit{\celsius})},
            ylabel={Solubilidad\vspace*{5pt}\\ \( \left(  \dfrac{\unit{g} \text{of solute}}{\qty{100}{\g} \text{ de agua}} \right) \)},
            ylabel style={
                at={(axis description cs:0.1,1.09)},
                anchor=south east,
                xshift=-15pt,
                rotate=-90,
                align=center,
                font=\small
            },
            xlabel near ticks,
            ]

% NaOH
            \addplot[draw, black, smooth, thick] coordinates {(0,15) (25,25) (60,40) (100,50) (160,71)};
            \node[font=\footnotesize, anchor=west] at (axis cs:140,60) {\ch{NaOH}};

% NaCl
            \addplot[blue, smooth, thick] coordinates {(0,35.7) (20,35.9) (40,36.2) (60,36.5) (80,36.8) (100,37) (160,38)};
            \node[blue, font=\footnotesize, anchor=south] at (axis cs:132,39) {\ch{NaCl}};

% KNO3
            \addplot[red, smooth, thick] coordinates {(0,13.3) (20,31.6) (40,63.9) (60,108) (77,160)};
            \node[red, font=\footnotesize, anchor=east] at (axis cs:72,145) {\ch{KNO3}};

% KCl
            \addplot[green!60!black, smooth, thick] coordinates {(0,28.0) (20,34.0) (40,40.5) (60,46.5) (80,52.0) (100,56.5) (160,74)};
            \node[green!60!black, font=\footnotesize, anchor=north west] at (axis cs:115,71) {\ch{KCl}};

% NH4Cl
            \addplot[orange, smooth, thick] coordinates {(0,29.0) (20,37.0) (40,45.0) (60,55.0) (80,66.0) (100,76.0) (160,105)};
            \node[orange, font=\footnotesize, anchor=south west] at (axis cs:111,89) {\ch{NH4Cl}};
        \end{axis}
    \end{tikzpicture}
    \caption{Solubility Curve}
    \label{fig:curva_solubilidad}
\end{figure}



